\clearpage
\pdfbookmark[1]{Coytier. 1813 März 6.}{anzeige_coytier_125_127}
\SourceLeaf{125}
\begin{center}
\par\noindent\hfill\rule{0.30\linewidth}{0.4pt}\hfill\mbox{}\par
Göttingische gelehrte Anzeigen. 1813 März 6.
\par\noindent\hfill\rule{0.30\linewidth}{0.4pt}\hfill\mbox{}\par
\end{center}

\textit{Recherches mathématiques ou diverses questions non résolues ou dont la solution} \textit{laisse quelque chose à désirer.} Par P. L. \textsc{Coytier.} \textit{Premier mémoire contenant} \textit{des observations générales sur les équations algébriques.} Paris 1813. \textit{Chez Eberhardt.} (16 Seiten in Quart.)\par

Wenn wir die wortreiche, schwülstige Vorrede recht verstehen, so soll der Zweck des gegenwärtigen Aufsatzes, und anderer, die ihm noch folgen sollen, dahin gehen, allerlei nach des Verfassers Meinung usurpirte Resultate der Analyse mit den Waffen der Synthese zu bekämpfen und umzustürzen. Aus dem Ganzen der Schrift sieht man aber, dass der Verfasser eigentlich mit den Wörtern Analyse und Synthese ihm eigenthümliche Begriffe verbindet, und bei jener sich ein blosses blindes Zeichenspiel, bei dieser eine reelle anschauliche Erkenntniss denkt. Ein solches Unternehmen verdient allerdings Lob. wenn sich ein solches blindes Zeichenspiel vorfindet, das sich den Namen Analyse anmasst; bei den meisten Versuchen dieser Art aber, diesseit und jenseit des Rheins gemacht, finden wir, dass die Blindheit nicht in dem Bekämpften, sondern in dem Bekämpfenden lag. Wohin gegenwärtige Schrift zu rechnen sei, möge man aus dem Inhalt, den wir kurz anzeigen wollen, beurtheilen.\par

Hauptsächlich ist die vorliegende Schrift gegen die Zerlegbarkeit der algebraischen Gleichungen in einfache Factoren gerichtet (schon in diesem Ausdrucke liegt eine Unrichtigkeit, die sich freilich auch manche andere Schriftsteller haben
\SourceLeaf{126}
zu Schulden kommen lassen). Der Verfasser nennt dieselbe auch öfters das Theorem von 1746, ohne Zweifel, weil \textsc{d’Alembert} um diese Zeit zuerst den Versuch eines strengen Beweises machte. Dass dieser Beweis von \textsc{d’Alembert}, eben so wie die Beweise von \textsc{Euler}, \textsc{Foncenex}, \textsc{Lagrange}, keineswegs befriedige, darüber haben wir schon vor 14 Jahren an einem andern Orte unser Urtheil erklärt, und eben so müssen wir freilich, unserer Ueberzeugung nach, von \textsc{Laplace}’s und \textsc{Lagrange}’s spätern Arbeiten über denselben Gegenstand urtheilen. Allein in diese Beweise selbst lässt sich unser Verfasser gar nicht ein: seine Bemerkungen sind ganz anderer Art, und nicht gegen die Beweise, sondern gegen den Lehrsatz selbst gerichtet. Er behauptet nemlich (um nur bei dem einfachsten Fall einer quadratischen Gleichung $xx-2ax+b=0$ stehen zu bleiben), $x-[a+\surd(aa-b)]$ sei gar kein einfacher Factor, weil $x-[a+\surd(aa-b)]=0$ gar keine Gleichung der ersten Ordnung, sondern eine wahre quadratische Gleichung sei. Man möge vor die Wurzelgrösse das Zeichen $+$, oder das Doppelzeichen $\pm$ schreiben, immer drücke sie beide Wurzeln zugleich aus. Jede Gleichung stelle eigentlich die Relation zwischen zwei veränderlichen und einer beständigen Lineargrösse dar, die Classification der Gleichungen und die Classification der Curven nach Ordnungen müsse aufs genaueste zusammenhangen, die Gleichung $x-[a+\surd(aa-b)]=0$ als identisch mit der Gleichung $xx-2ax+b=0$ betrachtet, und also erstere so gut, wie letztere, zur zweiten Ordnung gezählt werden. Diese Behauptungen machen den Inhalt der Schrift aus, und zeigen uns nichts, als die Verworrenheit der Begriffe des Verfassers. Die gemeine Algebra kennt gar keine veränderlichen Grössen, sondern bloss unbekannte und bekannte. Das Wurzelzeichen $\surd$ hat eigentlich in der mathematischen Zeichensprache eine doppelte Bedeutung (und dies ist allerdings eine kleine Unvollkommenheit); $\surd A$ soll entweder definirt werden, \textit{eine} Grösse, deren Quadrat $=A$, oder \textit{die} positive Grösse, deren Quadrat $=A$, insofern $A$ positiv ist. Es hängt von dem Analysten ab, wie er das Zeichen gebrauchen will, und ein denkender, vorsichtiger Analyst wird sich immer klar bewusst sein. und sich immer so ausdrücken, dass auch dem Leser kein Zweifel übrig bleibe, ob das Zeichen in unbestimmter oder in bestimmter Bedeutung gebraucht sei. \textit{Gleichungen} werden, wie wir schon oben andeuteten, gar nicht in Factoren zerlegt, sondern Functionen \textit{einer} veränderlichen Grösse. Also nicht die Gleichung $xx-2ax+b=0$, sondern die Function $xx-2ax+b$ wird in die Factoren $x-[a+\surd(aa-b)]$, $x-[a-\surd(aa-b)]$ zerlegt, und inso\SourceLeaf{127}%
fern man darin bloss $x$ als veränderlich betrachtet, \textit{nennt} man dieselben Factoren der ersten Ordnung. Immerhin mag der Verfasser, wenn er die Gleichung $xx-2ax+b=0$ als Ausdruck einer Curve, und $a$ oder $b$ als veränderlich betrachtet, den Factor $x-[a+\surd(aa-b)]$ einen Ausdruck der zweiten Ordnung nennen, er bestreitet dadurch keine \textit{Wahrheit}, sondern nur die Schicklichkeit einer \textit{Benennung} in einem Fall, worin man sie nicht gebraucht, und nie gebraucht hat.\par

Den meisten mathematischen Lesern dieser Blätter werden freilich diese Erörterungen überflüssig sein; wir glauben aber doch die Begriffe des Verfassers etwas entwickeln zu müssen, damit Niemand sich durch den Titel der Schrift verleiten lasse, neue Aufklärungen darin zu erwarten, wozu dem Verfasser die ersten Elemente zu fehlen scheinen.\par
\par\noindent\hfill\rule{0.30\linewidth}{0.4pt}\hfill\mbox{}\par

